
\chapter{Telepítés és használat}

A Magyar Ispell fejlesztés eredményei Linux, Windows és Mac OS X
operációs rendszereken is használhatóak olyan programokkal,
mint az OpenOffice.org, AbiWord, Hunspell, Emacs, vagy a magyar nyelv
támogatása szempontjából elavult Ispell.

\section{OpenOffice.org}

Az OpenOffice.org a legjobb
nyitott forráskódú irodai csomag, ezért a legtöbb felhasználó
az OpenOffice.org révén kerül kapcsolatba a Magyar Ispell
fejlesztés eredményeivel.

\subsection{Telepítés}

Az OpenOffice.org hivatalos magyar változata az FSF.hu Alapítvány
honlapjáról érhető el: \url{http://office.fsf.hu}.
Itt tölthető le Windows, és Linux platformra is, illetve
már letöltött formában számítástechnikai lapok
mellékletén is megtaláljuk.

A program magyar nyelvű telepítővel rendelkezik.
Kövessük a telepítő, vagy a honlap (illetve az újság)
útmutatását!

\subsection{Használat}

Az OpenOffice.org magyar változataiban alapértelmezett
nyelv a magyar, és rendszerint a beírás közbeni
ellenőrzés is be van kapcsolva. Ez utóbbi lehetőséget
a bal oldali eszköztáron található piros hullámvonallal aláhúzott \textsc{ABC} ikonnal
tudjuk ki- és újra bekapcsolni.

Az OpenOffice.org kis hibája, hogy a szövegszerkesztőben az
első betű beírása után töltődik be a magyar helyesírási
szótár. Ez lassabb gépeken másodperceket is igénybe vesz.
Mivel az ellenőrzés párhuzamosan folyik az írással, nem
kell megvárnunk a szótár betöltődését, nyugodtan elkezdhetünk gépelni.
(Igaz, a gépelt szöveg egy részét a szótár betöltődésének befejeződéséig
nem fogjuk látni.)

Ha beírás közben hibás szót írtunk be, a szó piros hullámvonallal
kerül aláhúzásra. Ha a jobb egérgombbal kattintunk a szón, a megjelenő
helyi menü tartalmazza a javaslatokat a hibás szóra.

Az ellenőrzést külön párbeszédablak segítségével is elvégezhetjük.
Ez az ablak az F7-es gomb lenyomásával hívható elő a leggyorsabban.
A később tárgyalt beállításokhoz itt egy helyen férhetünk hozzá,
így célszerű lehet megismerkedni a használatával.

\subsection{A szöveg nyelvének kiválasztása}

Előfordul, hogy olyan dokumentummal találkozunk, ahol a szöveg,
vagy annak egy része hibásként jelenik meg, bár helyes magyar
szavakat tartalmaz. Ennek oka, hogy a szöveg -- bár magyar --
más nyelvűnek lett beállítva.

Jelöljük ki a szöveget, és válasszuk ki a Formátum$\rightarrow$Karakter
párbeszédablakban a Betűtípus lapot! Válasszuk ki a Nyelv
listából a magyart, és kattintsunk a \textsc{OK} gombon!

Ha egy szó, vagy bekezdés nyelvét szeretnénk átállítani,
az OpenOffice.org rendelkezik egy kényelmes lehetőséggel:
kattintsunk egy hibásnak jelzett szón a jobb gombbal. Az OpenOffice.org
a rendelkezésre álló helyesírási szótárak alapján megpróbálja megállapítani,
hogy milyen nyelvű a kérdéses szó. Ha ez sikerül, a felnyíló helyi
menü alján két új menüpontot is láthatunk. A példánknál maradva, ,,A szó: magyar''
menüponttal a kiválasztott szó, a ,,bekezdés: magyar'' menüponttal a teljes bekezdés nyelvét
állíthatjuk egy kattintással magyarra.

\subsection{Nagybetűs és számokat tartalmazó szavak ellenőrzése}

Az Eszközök$\rightarrow$Beállítások párbeszédablakban a
Nyelvi beállítások$\rightarrow$Írástámogatás lapon egyéb
beállításokat is megtehetünk. A nagybetűs és számokat tartalmazó
szavak ellenőrzését a Beállítások dobozban lehet kérni.
Ezt igényesebb dokumentumok ellenőrzése esetén érdemes megtenni.

(A korábbi változatok
a nagybetűs szavakat helyességtől függetlenül hibásnak jelezték,
ha kértük az ilyen szavak ellenőrzését. Ez a probléma
az OpenOffice.org 1.02 magyar változatával megszűnt.)

\subsection{Alapértelmezett nyelv beállítása}

Az Eszközök$\rightarrow$Beállítások párbeszédablakban a
Nyelvi beállítások$\rightarrow$Nyelvek lapon állíthatjuk be az OpenOffice.org
dokumentumok alapértelmezett nyelvét. Jelöljük be a Csak az aktuális
dokumentumban jelölőnégyzetet, ha csak az éppen szerkesztett
dokumentum alapértelmezett nyelvét szeretnénk megváltoztatni!

\subsection{Automatikus elválasztás beállítása}

A helyesírás része a helyes elválasztás. Ennek automatikus változatát
az Eszközök$\rightarrow$Beállítások párbeszédablakban a
Nyelvi beállítások$\rightarrow$Írástámogatás lapon választhatjuk ki.
Kattintsunk a Szerkesztés gombon, a megjelenő párbeszédablakban válasszuk ki
a magyar nyelvet, majd az ALTLinux LibHnj Hyphenator jelölőnégyzetet
válasszuk ki.

Az OpenOffice.org magyar változatai között akadnak olyanok is, ahol az
előző beállítás nem működik (könnyen ellenőrizhető: kattintsunk
újra a Szerkesztés gombon, ha az előző jelölőnégyzet újra
nincs kiválasztva, a beállítás nem lehetséges). Ilyenkor ha mégis
elválasztást szeretnénk, jelöljük ki az egész dokumentumot,
és a Formátum$\rightarrow$Bekezdés
párbeszédablakban válasszuk ki a Szövegfolyam lapot, és jelöljük
be az Elválasztás szakaszban az Automatikusan jelölőnégyzetet!

\subsection{Helyesírási állományok frissítése}

A helyesírási állományokat (\verb+hu_HU.dic+, \verb+hu_HU.aff+,
\verb+hyph_hu.dic+) a Magyar Ispell honlapról
(\url{http://magyarispell.sf.net}) tölthetjük le, és az OpenOffice.org
\texttt{share/dict/ooo} (régebben \texttt{user/wordbook}) könyvtárába
másoljuk át kibontás után. Az OpenOffice.org újraindítás
után már az új állományokat fogja használni.

%%%%%%%%%%%%%%%%%%%%%%%%%%%%%%%%%%%%%%%%%%%%%%%%%%%%%%%%%%%%%%%%%

\subsection{Az OpenOffice.org nem magyar változatai}

\subsubsection{Különbségek}

Az OpenOffice.org nem magyar változatai nem tartalmazzák
az OpenOffice.org Myspell nevű helyesírás-ellenőrzőjének
magyar bővítését (amelyen tulajdonképpen a Hunspell program is
alapszik), emiatt sok hibát követnek el a helyesírás
ellenőrzésében, különösen az összetett szavak
kezelése terén. Például nincs korlátozva az összetett
szavak hossza; nem ismeri fel az olyan tipikus hibákat,
mint a \emph{színtű}, \emph{karvaj}, \emph{rovarirtószer}; sok
helyes szót sem fogad el, mint például
\emph{kéthavi}, \emph{Sydneyvel}; nem kezeli a kötőjeles
szavakat, mint a \emph{MÁV-ot}, \emph{Voltaire-rel} stb.

Könnyen meggyőződhetünk róla, hogy az OpenOffice.org nem magyar
változatával van dolgunk: írjuk be a szövegszerkesztőbe, hogy
\emph{helyesírásellenőrző}. A nem magyar változatok helyesnek
fogadják el  a szót, ha egyébként a magyar helyesírási szótár
lett kiválasztva az Eszközök$\rightarrow$Beállítások párbeszédablakban a
Nyelvi beállítások$\rightarrow$Írástámogatás lapon. A
magyar változatok hibásnak jelzik a szót, és a
\emph{helyesírás-ellenőrző} alakot javasolják helyette.

\subsubsection{Helyesírási állományok telepítése}

Ha mindezek ellenére a nem magyar változat használata mellett döntünk,
szükségünk lehet a helyesírási állományok (\verb+hu_HU.dic+, \verb+hu_HU.aff+,
\verb+hyph_hu.dic+) kézi telepítésére, ha ezeket nem tartalmazza az OpenOffice
változat. Az említett állományokat a Magyar Ispell honlapról
(\url{http://magyarispell.sf.net}) tölthetjük le, és az OpenOffice.org
megfelelő könyvtárában kell elhelyeznünk.

A könyvtár helye az újabb változatokban a \texttt{share/dict/ooo}.
Az itt található \emph{dictionary.lst} állományba kell még bejegyeznünk
a helyesírási állományokat a

\begin{verbatim}
DICT hu HU hu_HU
HYPH hu HU hyph_hu
\end{verbatim}

sorokkal. Az OpenOffice.org indítása után az
Eszközök$\rightarrow$Beállítások párbeszédablakban a
Nyelvi beállítások$\rightarrow$Írástámogatás lapon válasszuk ki a Szerkesztés
gombot, és a megjelenő lapon a magyar nyelvhez kapcsoljuk be az
OpenOffice MySpell SpellCheckert és az ALTLinux LibHnj Hyphenatort,
és az alapértelmezett nyelvet is módosítsuk a korábban ismertetett
módon, a Nyelvi beállítások$\rightarrow$Nyelvek lapon.

%%%%%%%%%%%%%%%%%%%%%%%%%%%%%%%%%%%%%%%%%%%%%%%%%%%%%%%%%%%%%%%%%

\section{A Hunspell telepítése és használata}

\subsection{Linux és egyéb Unix-típusú rendszerek alatt}

A Hunspell már része egyes Linux terjesztéseknek,
például az UHU-Linuxnak (az 1.0-s változatban még
Mispell néven). A legfrissebb változatokat
forráskódból fordíthatjuk le.

\subsubsection{Fordítás}

Töltsük le a \url{http://magyarispell.sf.net} oldalról a
legfrissebb Hunspell csomagot.

Telepítsük a következő GNU, vagy ezzel egyenértékű fejlesztői
csomagokat: \texttt{glibc-devel}, \texttt{make}, \texttt{libncurses5-devel},
\texttt{readline-devel}, \texttt{gettext-devel}.

Csomagoljuk ki, és fordítsuk le a Hunspell forrást:

\begin{verbatim}
tar xvzf hunspell.tar.gz
cd hunspell
make
make install
\end{verbatim}

Amennyiben a telepítési útvonalak nem megfelelőek, a
Makefile állományban pontosíthatunk ezeken.

Ellenőrizzük a program működését a

\begin{verbatim}
hunspell latin12.html
\end{verbatim}

vagy a

\begin{verbatim}
hunspell ooomagyarispellteszt.txt
\end{verbatim}

parancs kiadásával.

\subsubsection{Fordítás menüs felület nélkül}

A \texttt{hunspell.cxx} állományban kapcsoljuk ki a
\textsc{CURSES} makrót, megjegyzésbe téve, vagy törölve a

\begin{verbatim}
#define CURSES
\end{verbatim}

sort. Ez esetben a menüs felület nem fog működni, csak
a parancssori. Fordítsuk és telepítsük a programot az előző
szakaszban ismertetett módon.

\subsubsection{A Hunspell használata}

A program és szótárállománya formátumának leírását a két
feltelepített kézikönyvoldal (Hunspell(1), Hunspell(4))
tartalmazza.

A telepítés során kiadott \texttt{make install} parancs az
Ispell parancshoz is a Hunspell programot
állítja be. Az eredeti Ispell \texttt{ispell.orig} néven indítható
a továbbiakban.
 
A személyes szótárat a \texttt{\$HOME/.hunspell\_default}, vagy
ha nem az alapértelmezett magyar szótárakat használjuk, akkor a
\texttt{\$HOME/.hunspell\_[nyelv]} fájlba írja a program.

A saját szótár később kézzel is átszerkeszthető, érdekes
tulajdonsága, hogy ragozható szavakat is felvehetünk ide.

A Hunspell az Ispellhez hasonlóan együttműködik 
az Emacs, LyX, KWord programokkal. A KWordhöz való
beállításban Szántó Tamás leírása segít. A leírás a
\texttt{http://magyarispell.sf.net} oldalon található.

Ha a Hunspell programot állománynév paraméterekkel
hívjuk meg, sorban megnyitja az állományokat, és a hibás
szavakat egy menüs felület segítségével felkínálja
javításra.

A kiterjesztés alapján felismeri a szöveges,
\TeX{}, Nroff/Troff és HTML állományokat, de kapcsolókkal
is befolyásolhatjuk ezt: \emph{-t} (\TeX{}),
\emph{-H} (HTML), \emph{-n} (Nroff/Troff).

A felület további ismertetését a függelékben található
Hunspell kézikönyvoldal tartalmazza. Itt csak a
használható billentyűkombinációk leírását idézzük.
Zárójelben megadjuk a billentyűparancsok angol
megfelelőit is.

\begin{itemize}
\item  \emph{Szóköz}: Továbblépés javítás nélkül a következő hibás szóra.

\item  \emph{Számok}: A megfelelő javaslat kiválasztása.

\item  \emph{c} (r):      Csere. A hibás szó cseréje a megadott  szóra,  vagy
              szavakra.     Az    Escape    billentyű    kétszeri
              lenyomásával megszakíthatjuk a a cserét.

       \item  \emph{j} (a):            A    szó    elfogadása    a    program    futásának
              befejeződéséig.

\item  \emph{f} (i):            A szó felvétele a saját szótárba.

\item  \emph{k} (u):            A szó kisbetűsként való felvétele a saját szótárba.

\item  \emph{t} (--):            Ragozott tőszó felvétele. A ragozás  a  másodikként
              megadott,  mintául  szolgáló  szó alapján történik.
              Ha  ez  a  szó  nincs   felvéve   a   szótárba,   a
              tőszófelvételt elölről kell kezdeni.

              Példa:   gyakran   hiányoznak   az   idegen   és  a
              mozaikszavak. Vegyük fel a hiányzó OTP  szót  a  tv
              köznévi  mozaikszó  toldalékaival!   Lenyomjuk  a t
              billentyűt, beírjuk, hogy OTP, Enter, beírjuk, hogy
              tv,  és  Enter.  A  tv  szó  ugyanis már szerepel a
              szótárban, és ragozása  mondja  meg  a  programnak,
              hogy  az  OTP  szót milyen toldalékokkal szeretnénk
              látni, pl.  OTP-hez, OTP-met stb. a tv-hez,  tv-met
              mintájára.   A  -val/-vel,  -vá/-vé  toldalékokat a
              kötőjellel kapcsolt toldalékú szavak esetében külön
              kell  felvenni, de még így is sokkal kényelmesebb a
              speciális tőfelvétel, mint minden (akár több  száz)
              toldalékos alakot külön-külön felvenni.

              Egy  kis  segítség  a  minta  kiválasztásához: első
              menetben a hangrend számít. Pl.  mozaikszavaknál  a
              mély hangrendű új szavak felvételénél használhatjuk
              a MÁV szót példaként, ajakréses  magas  hangrendnél
              pedig   az  említett  tv-t.  Ajakkerekítéses  magas
              hangrend esetében a HÖK szót használhatjuk (amíg át
              nem lesz javítva a helyes hök-re). Ha a hiányzó szó
              nem mozaikszó, hanem például  tulajdonnév,  a  Pál,
              Péter,   Örs   szavakat  adhatjuk  meg  példának  a
              hangrendtől függően. Ha meg egyszerű köznévről  van
              szó,  akkor  használjuk  az  ok,  ék,  ük szavakat.
              Igéknél áll, vél, ül, ikes  igéknél  ázik,  vérzik,
              nőzik  lehet  a  példa  (remélhetőleg  a  következő
              változatokban egyszerűsödik ez a tőfelvétel).

              Az   Escape   billentyű    kétszeri    lenyomásával
              megszakíthatjuk ezt a műveletet.
	      
Az itt felvett szavak automatikusan a saját
szótárba kerülnek. Ha ezt innen törölni
szeretnénk, a programból kilépve
a saját szótár kézi átszerkesztésével tehetjük meg.

\item  \emph{m} (q):            Kilépés  a  változtatások mentése nélkül. A program
              változtatás megléte esetén megerősítést kér.

\item  \emph{v} (x):            Az   állomány   ellenőrzésének    megszakítása    a
              változtatások    mentésével.   Ha   van   következő
              állomány, akkor a program annak ellenőrzését  kezdi
              el.

\item  \emph{?} (?):            Segítség.  Hasonló  rövid  leírás  megjelenítése  a
              billentyűparancsokról.
\end{itemize}

\subsection{Windows alatt}

Töltsük le a \url{http://magyarispell.sf.net} oldalról a
Windowsra készített csomagot, és csomagoljuk ki
a \texttt{C:\\} meghajtó gyökerében a WinZip programmal.

A kicsomagolás után a futtatható
program (\texttt{Hunspell}) és a szótárállományok a
\texttt{C:\\Myspell} könyvtárba kerülnek.

Nyissunk meg egy Parancssor ablakot, majd adjuk ki a
\texttt{C:\\Myspell\\Hunspell} parancsot. A beírt szavakat,
illetve mondatokat a Hunspell feldolgozza, és a hibás
szavakat \&, vagy \# jellel kezdve írja ki. (A kimeneti
formátum részletes ismertetését a függelékben közölt
kézikönyvlapok tartalmazzák.)

Ha egy meglévő állományban szeretnénk a hibás szavakat
megkeresni, a \emph{-l} kapcsolót, és esetleg az
állomány típusát adjuk meg, az állományt pedig a
standard bemeneten keresztül adjuk a Hunspell programnak.

\begin{verbatim}
hunspell -l -H <próba.html
\end{verbatim}

A példában a \texttt{próba.html} állomány hibás szavait
irattuk ki a képernyőre. Ha a kimenetet el szeretnénk
menteni, irányítsuk fájlba:

\begin{verbatim}
hunspell -l -H <próba.html >hibás_szavak.txt
\end{verbatim}

A Windows alá letölthető futtatható állomány nem tartalmazza
a Linux alatt alapértelmezésként működő menüs felületet, amely
lehetővé teszi a szöveges, \textsc{HTML}, \LaTeX{} és
Groff formátumú állományok gyors javítását.\footnote{A Cygwin
környezet lehetővé teszi a menüs rendszer fordítását.
A menüs rendszer Ncurses, Readline,
és Gettext függvénykönyvtárakat használ, amelyeket az Visual C++
-- szemben a Cygwinnel -- alapból nem tartalmaz.
A Cygwin címe: \url{http://www.cygwin.com}.
}


\subsection{Hunspell Mac OS X alatt}

Fazekas Mihály (\url{boga@mac.com}) a
Hunspell 0.94-es változatát egy Cocoa felülettel bővítette,
ezzel lehetővé téve a helyesírás-ellenőrző rendszerszintű
használatát Mac OS X alatt. Letölthető csomagok, és részletes telepítési
útmutató található a \url{http://myspellx.sf.com} címen.

A Hunspell program lefordítható Mac OS X alá is. Telepítsük
a GNU fejlesztőeszközöket (glibc-devel, ncurses-devel,
readline-devel, gettext-devel), vagy a \texttt{hunspell.cxx}
állományban kapcsoljuk ki a \textsc{CURSES} makrót, megjegyzésbe
téve, vagy törölve a

\begin{verbatim}
#define CURSES
\end{verbatim}

sort. Ez utóbbi esetben a menüs felület nem fog működni, csak
a parancssori.

%%%%%%%%%%%%%%%%%%%%%%%%%%%%%%%%%%%%%%%%%%%%%%%%%%%%%%%%%%%%%%%%%

\section{Emacs}

Az \texttt{Emacs}, illetve a \texttt{Xemacs} program mind Windows,
mind Linux alatt képes a Hunspell programot használni.

Ismeretlen okból csak a Xemacs kezeli jól
a kötőjeleket tartalmazó szavakat, így ennek használata
javasolt az Emacs helyett.

\subsection{Az Emacs telepítése Linuxra}

Linux esetében a terjesztések tartalmazzák mindkét programot.
Az alapcsomag az \texttt{emacs}. Ha az \texttt{emacs-X11}
csomagot is telepítjük, az X grafikus felületen egy bővebb
változat indul el, ami menüsorral és eszköztárral is rendelkezik.
Ez utóbbival jelentős részben megegyezik a \texttt{xemacs} csomagban található
\texttt{Xemacs}.

\subsection{Az Emacs telepítése Windowsra}

Windowsra több helyen is találhatunk bináris csomagot, például az
\url{ftp://ftp.gnu.org/gnu/windows/emacs/latest/} címen.
Általános esetben az \texttt{emacs-21.2-bin-i386.tar.gz} csomagra van
szükség, melyet az
\url{ftp://ftp.gnu.org/gnu/windows/emacs/utilities/i386/} címen
található programokkal  vagy a WinZIP programmal is kibonthatunk.

Egy másik lelőhely a CTAN archívum, például:
\url{ftp://ftp.dante.de/tex-archive/systems/win32/fptex/0.5/support/emacs-21.1.1-win32.zip}.

\subsection{A Xemacs telepítése Windowsra}

A \url{http://www.xemacs.org} címről tölthető le a Xemacs
windowsos változata.

\subsection{A Hunspell telepítése az Emacs alá}

Az Emacs és Xemacs programok a háttérben indított Hunspell-lel
együttműködve képesek a helyesírás-ellenőrzésre.

A Hunspell-nek a keresési útvonalon kell elhelyezkednie, méghozzá \texttt{ispell}
néven. A Linuxra történő telepítés során ez megtörénik.
Windows alatt készítsünk a \texttt{C:\\Myspell\\Hunspell} programról
egy másolatot \texttt{Ispell} néven valamelyik rendszerkönyvtárba
(például \texttt{C:\\Windows}, \texttt{C:\\WinNT}).

Windows, és az Emacs X-en futó változatai (Emacs-X11, Xemacs)
esetében további beállításra nincs szükség. Az \texttt{M-x\ ispell}
parancs kiadása után az Emacs elindítja a Hunspell programot,
és ellenőrzi a segítségével az éppen szerkesztett szöveget.

A Xemacs esetében a \texttt{M-x\ ispell-buffer} parancsot
adjuk ki elsőnek, majd a következő ellenőrzések
során már használható \texttt{M-x\ ispell} is.

\subsection{A Hunspell telepítése konzolon indított Emacshoz}

Az Emacs (tehát nem Xemacs) konzolon való használathoz szükségünk
lehet az \texttt{ispell.el}
módosított változatára, bár inkább javasolt helyette a
Xemacs konzolon történő indítása.

Az ispell.el állomány a \url{http://magyarispell.sf.net/ispell.el} címen
található meg. A letöltött és kibontott állományt másoljuk abba a könyvtárba, ahol
a lefordított \texttt{ispell.elc} található (például
\texttt{/usr/share/emacs-21.1/lisp/textmodes/} könyvtár, vagy régebben
\texttt{/usr/share/emacs/20.7/lisp/}).

Ezután töröljük az ebben a könyvtárban található \texttt{ispell.elc} állományt!

A telepített Ispell.el \texttt{magyar} nevű szótárállományokat
keres, ezért győzödjünk meg arról, hogy a default.aff, default.dic
állományok ilyen néven is elérhetők (Windows alatt a
\texttt{C:\\Myspell}, Linux alatt a \texttt{/usr/share/myspell} könyvtárakban).
Ha nem, akkor készítsünk másolatot (vagy kötést) ezekről \texttt{magyar.aff}
és \texttt{magyar.dic} néven ugyanabban a könyvtárban.

Adjuk ki az \texttt{M-x} \texttt{ispell-change-dictionary} parancsot,
és válasszuk ki a magyar szó beírásával a magyar Ispell szótármodult.

Ha magyar az alapbeállításunk, természetesen az Emacs indító
konfigurációjába (vagyis a \verb+~/.emacs+ fájlba) is beleírhatjuk:

\begin{verbatim}
(ispell-change-dictionary "magyar")
(ispell)
(ispell-kill-ispell)
\end{verbatim}

A folyamatos ellenőrzés alatt használható billentyűparancsok
megegyeznek a Hunspell angol billentyűparancsaival.

A saját szótárba felvett szavak
a \texttt{\~/.hunspell\_default}, a módosított ispell.el használata esetén
a \texttt{\~/.hunspell\_magyar} állományba kerülnek.

Minthogy a Spell-menü tartalmának alapbeállítása az emacs programba
bele van fordítva, indításkor az nem tartalmazza a magyar lehetőséget,
viszont az ispell legalább egyszeri indítása után már igen, tehát
ekkor már a menü is használható a parancssori utasítások helyett: a
20.\ verzió esetében Edit$\rightarrow$Spell, a 21.\ verzió esetében
pedig Tools$\rightarrow$Spell Checking.

\subsection{Az Emacs használata}

A teljes szöveg ellenőrzésére szolgál az \texttt{M-x ispell} parancs.

Az itt használható billentyűparancsok
megegyeznek a Hunspell angol billentyűparancsaival.

Egy szót az \texttt{M-x} \texttt{ispell-word} paranccsal 
ellenőriztethetünk. Ehhez létezik gyorsbillentyű is: \texttt{M-\$}.

Az Emacs \emph{flyspell} üzemmódja lehetőséget nyújt a beírás
közbeni helyesírás-ellenőrzésre, de a magyar beállításokkal
nem működik tökéletesen együtt. Az üzemmódba az
\texttt{M-x} \texttt{flyspell-mode} paranccsal léphetünk be.

Ezután a hibás szavak grafikus felületen
piros színűvé és aláhúzottá válnak. Ha a 3.\ egérgombbal (vagy a
két szélsővel egyszerre) kattintunk a szón, felnyílik egy
menü, amellyel a javasolt cseréket fogadhatjuk el, illetve továbbléphetünk,
felvehetjük a hibás szót az elfogadott szavak közé az ellenőrzés idejére,
illetve ,,végérvényesen''.

%%%%%%%%%%%%%%%%%%%%%%%%%%%%%%%%%%%%%%%%%%%%%%%%%%%%%%%%%%%%%%%%%

\section{Az Ispell telepítése és használata}

Az Ispell a Hunspell képességeit tekintve elavult, és
magyar nyelvű szövegek helyesírás-ellenőrzésére csak
végszükség esetén, a Hunspell hiányában javasolt. A különbségeket a
függelékben található Hunspell kézikönyvoldal foglalja
össze.

\subsection{Telepítés Linuxra}

Az Ispell része a legnagyobb Linux terjesztéseknek,
és egy átlagos Linux telepítésnél az Ispell telepítésére
is sor kerül automatikusan.

Linux alatt az Ispell szótárállományai a \texttt{/usr/lib/ispell}
állományban kerülnek elhelyezésre. A magyar Ispell szótár,
vagyis a (\texttt{magyar.hash}) állomány is ide kerül.

\subsection{Telepítés Linuxra, forráskódból}
Az ispell-honlapról
(\url{http://fmg-www.cs.ucla.edu/fmg-members/geoff/tars/}) letölthető
forráscsomag a benne lévő \texttt{README} fájlban leírtak szerint
(make all ; su ; make install) lefordítható.  A local.h fájl linux
esetében javasolható módosításai a
\begin{verbatim}
#define NO8BIT    /* Remove this if you use ISO character sets */
\end{verbatim}
\noindent sor törlése és a 
\begin{verbatim}
#define USG       /* Define on System V or if term.c won't compile */
#define LIBDIR    "/usr/lib/ispell"
\end{verbatim}
\noindent sorok eredetijeinek átírása a közöltekké.

Némely rendszerek esetében szükség lehet a \texttt{/usr/lib/} könyvtárban a 
\verb+# ln -s libtermcap.so.2 libtermcap.so+ kötés létrehozására.

Ezután már csak make és -- rootként -- make install, mely parancs a
\texttt{/usr/lib/ispell/}-be teszi az angol adatfájlokat és 
/usr/local/bin/-be a programokat.

\subsection{Windows\label{w32-ispell}}

A Windowsra történő telepítéshez Geoff Kuenning Ispell honlapján
(\url{http://fmg-www.cs.ucla.edu/fmg-members/geoff/ispell-winnt.html})
több lehetőséget is találunk.

Windows alatt két fő módja van az Ispell telepítésének
és használatának. Az elsőben a teljes Cygwin környezetet
telepítjük a Windows NT, 2000, illetve XP operációs rendszerünkre.
Ez gyakorlatilag teljes Unix környezetté egészíti ki
a Windowst, és utána telepítünk egy Cygwin környezet alá
fordított Ispell változatot.

A másik módszer azok számára, akik a szokásos Windows környezetet
előnyben részesítik, egyszerűbb: ez a bináris csomag formájában
letölthető WIN32-re fordított Ispell változat Cygwin nélküli
használata.

Mindkét módszer esetében az
\url{ftp://ftp.franken.de/pub/win32/develop/gnuwin32/cygwin/%
porters/Humblet_Pierre_A/V1.1/} címen található verziót javasoljuk. A
telepítést az \texttt{ispell-3.2.06-cygwin-1.3-bin.README} nevű fájl
ismerteti. Az egyszerűbb esetben a csomagból kizárólag a
\texttt{usr/local/bin/ispell.exe} fájlra van szükség, melyet a
PATH-ban található valamelyik könyvtárba kell elhelyezni.  A
szótárfájlok (lásd \ref{MI-telep}) a \texttt{DICTDIR} környezeti
változó által meghatározható helyre teendők.  (A .tar.gz végződésű
fájl kicsomagolásához a Cygwin környezeten kívül pl.\ a djtarnt.exe
program, vagy akár a népszerű WinZIP újabb változatai használhatók.)

Ha tehát a \verb+C:\AUTOEXEC.BAT+ fájlban a
\begin{verbatim}
SET PATH=C:\WINDOWS;C:\WINDOWS\COMMAND;C:\BIN
SET DICTDIR=C:\ISPELL
\end{verbatim}
\noindent sorok találhatók, akkor az ispell.exe fájlt a \verb+c:\bin+, 
a magyar.hash és a magyar.aff fájlt a \verb+c:\ispell+ alkönyvtárba
kell helyezni.

Szintén a PATH-ban kell elhelyezni a cygwin1.dll nevű fájlt.

Egy másik bináris csomag formájában letölthető WIN32-re fordított
Ispell változat található a
\url{ftp://ftp.tue.nl/pub/tex/GB85/ispell-dutch96/ispellw32.zip}
címen.

A letöltés után csomagoljuk ki a ZIP állományt a
\verb|C:\| meghajtó gyökerében, és adjuk a
keresési útvonalunkhoz a \verb|C:\Ispell\bin|
könyvtárat. Ehhez nyissuk meg a 
Start menü$\rightarrow$Beállítások$\rightarrow$Rendszer
párbeszédablakot, és a Speciális lapon
a középső részben (Környezet)
keressük meg a keresési útvonalakat rögzítő környezeti
változót (PATH), és írjuk a végére
pontosvesszővel elválasztva a
\verb|C:\Ispell\bin| útvonalat.

Sajnos ez a dán Ispell modulhoz készült csomag
még kis számú ragozási osztállyal
lett lefordítva, ezért ha a legfrissebb
Magyar Ispell szótármodult szeretnénk használni vele,
újra kell fordítanunk az Ispellt. Ehhez
az EMX+gcc csomagokra van szükség.

Egy harmadik win32-változat a CTAN:fptex
(\url{ftp://ftp.dante.de/tex-archive/systems/win32/}) disztribúcióban
található.

\subsection{A \texttt{magyar.hash} állomány telepítése és használata}

A \texttt{magyar.hash} állomány, vagyis a magyar Ispell szótár része több Linux
terjesztésnek is (különböző neveken: \texttt{ispell-hu}, \texttt{ihungarian} stb.).

A \url{http://magyarispell.sf.net/} oldalról töltsük
le a legfrissebb változatát a Magyar Ispell csomagnak
(\texttt{magyarispell-XXX-tar.gz}), majd adjuk ki a
következő parancsokat:

\begin{verbatim}
tar xvzf magyarispell-XXX.tar.gz
cd magyarispell-XXX
make ispell
cp tmp/magyar4ispell.hash /usr/lib/ispell/magyar.hash
\end{verbatim}

Windows esetében a \texttt{magyar.hash} állomány a
\ref{w32-ispell} pontban leírt helyre kell tenni.

Ha a magyart alapértelmezett nyelvként szeretnénk látni,
akkor gondoskodjunk arról, hogy a szótárállományokat
tartalmazó könyvtárban a \texttt{default.hash}
a \texttt{magyar.hash}-ra mutasson.

\begin{verbatim}
rm /usr/lib/ispell/default.hash
ln -s /usr/lib/ispell/magyar.hash /usr/lib/ispell/default.hash
\end{verbatim}

\subsection{Az Ispell változat megfelelősége}

Ahogy már említésre került, az Ispell számos fordítás során
beállított paramétertől függően végzi tevékenységét.
Különösen régebbi Linux terjesztésekben, vagy Windows alatt
meg kell győződnünk arról, hogy az Ispell képes
helyesen lefordítani és kezelni
a Magyar Ispell forrásállományait.

Adjuk ki Linux/Unix alatt a következő parancsot:

\begin{verbatim}
ispell -vv | grep BIT
\end{verbatim}

Ha a bináris Ispell megfelel a céljainknak, akkor a következő
két sort kell látnunk:

\begin{verbatim}
       MASKBITS = 64
       !NO8BIT (8BIT)
\end{verbatim}

Az első sor jelzi, hogy elegendő számú ragozási osztályt kezel az Ispell
(a régebbi Ispell változatoknál a 32 volt az alapértelmezett).

A második sor jelzi, hogy a szótárban található szavak
tartalmazhatnak ékezetes betűket is.

Windows alatt a következő parancsot adjuk ki a
parancssorban:

\begin{verbatim}
ispell -vv | more
\end{verbatim}

A listában keressük meg a fenti két sort.
Ha a MASKBITS értéke 32, vagy a NO8BIT értéke igaz,
akkor a Magyar Ispell használatához szükség van
az Ispell újrafordítására.

\subsection{Az Ispell használata}

\subsubsection{Szöveges állományok}

A következő paranccsal kérhetjük szövegfájlok ellenőrzését:

\begin{verbatim}
ispell -d magyar fájl(ok)
\end{verbatim}

Ebben az esetben megjelenik egy kényelmes és gyors menü,
amely sorban felmutatja a hibásnak vélt szavakat.

Ha van javítási javaslat, akkor a javaslat előtt látható
számjegyek leütésével fogadjuk el a javaslatot.

A legfontosabb billentyűparancsok a következőek:

\begin{itemize}
\item szóköz: Átugrik a szón, azt változatlanul hagyva.
\item számok: Választás a javasolt cserék közül.
\item \emph{r}: Csere az általunk megadott szóra.
\item \emph{a}: Elfogadjuk a szót az ellenőrzés befejeződéséig.
\item \emph{i}: Felvesszük a saját szótárunkba. (Ez a \verb+~/.ispell_magyar+, 
vagy ha létezik az aktuális könyvtárban az \texttt{.ispell\_magyar}
állomány, akkor az.  Később -- ha elhagytuk az Ispell programot -- ez
akár kézzel is módosítható.)
\item \emph{u}: Felvesszük a saját szótárunkba, de kisbetűsre
alakítjuk előtte.
\item \emph{x}: A következő fájlra ugrunk. A félig ellenőrzött 
állománynak a fennmaradó része nem lesz kijavítva.
\item \emph{l}: A rendszerszótárban végezhetünk kereséseket.
A \texttt{*} jel segítségével mintázatot is megadhatunk, hasonlóan
a fájlrendszerben megadott útvonalakhoz: pl. az \texttt{L*}
mintázat minden L betűvel kezdődő szót kilistáz a
\texttt{/usr/dict/words} állományból. A rendszerszótár angol nyelvű.
\item \emph{q}: Kilépés mentés nélkül, ha a ,,q'' után ,,y''-t
nyomunk megerősítő kérdésre. A vizsgált állomány javításai elvesznek.
\end{itemize}

A Hunspell program hasonló felülettel rendelkezik Linux alatt,
de több újdonsággal is szolgál a magyar nyelv számára.
Leírást a Hunspell(1) kézikönyv oldalon találunk.

\subsubsection{\TeX{} állományok}

A következő paranccsal kérhetjük \TeX{}, illetve \LaTeX{} állományok
ellenőrzését:

\begin{verbatim}
ispell -t -d magyar fájl(ok)
\end{verbatim}

Amennyiben az állomány(ok) kiterjesztése .tex, a -t kapcsolót
nem szükséges megadni.

A -t kapcsoló hatására az Ispell nem törődik
a fordított perjellel kezdődő, ékezetes betűket nem tartalmazó
szavakkal (vagyis a TeX parancsokkal). Bizonyos parancsokat
felismer, és az utána következő egy, vagy két argumentummal
sem foglalkozik. A matematikai környezetben lévő szavakat
nem ellenőrzi az Ispell.

\subsubsection{HTML állományok}

A következő paranccsal kérhetjük HTML állományok
ellenőrzését:

\begin{verbatim}
ispell -H magyar fájl(ok)
\end{verbatim}

